 \setlength{\headheight}{14.49998pt}
 
\chapter{Introduzione}
Le abitazioni intelligenti mettono a disposizione un numero crescente di dispositivi, dati e possibilit\`a di controllo. Perch\'e tali possibilit\`a siano davvero accessibili, l'utente deve poter tradurre le proprie abitudini e preferenze in automazioni comprensibili, verificabili e coerenti con i limiti della casa reale. Questo elaborato affronta tale esigenza nel contesto di un Gemello Digitale per smart home, concentrandosi sull'integrazione di funzionalit\`a di \textit{Trigger-Action Programming} nell'interfaccia web del sistema.

Il lavoro parte dall'analisi delle automazioni domestiche e delle modalit\`a con cui possono essere create attraverso interfacce visuali, conversazionali e multimodali. Su questa base descrive l'architettura del Gemello Digitale e propone un flusso visuale che guida l'utente nella definizione di trigger, condizioni e azioni, permettendo di distinguere la composizione della regola dalla sua simulazione e dal successivo salvataggio.

Il Capitolo~1 presenta contesto, obiettivi, metodo e delimitazione del lavoro. Il Capitolo~2 discute lo stato dell'arte; il Capitolo~3 descrive l'architettura del sistema; il Capitolo~4 analizza l'interfaccia di partenza e illustra il \textit{builder} TAP progettato e integrato durante la tesi. Le conclusioni sintetizzano il contributo e le possibili evoluzioni.
